\documentclass{article}
\input{../../../../lib.sty}

\title{\texttt{NoisePrivacyMap<L2Distance<RBig>, PrivacyCurveDP> for ZExpFamily<2>}}
\author{\MichaelShoemateAuthor}
\date{}

\begin{document}

\maketitle

\contrib
Proves soundness of the implementation of
\rustdoc{measurements/noise/trait}{NoisePrivacyMap} for
\texttt{ZExpFamily<2>} with output measure
\rustdoc{measures/struct}{PrivacyCurveDP}.

\section{Hoare Triple}
\subsection*{Precondition}
\subsubsection*{Compiler-verified}
\texttt{NoisePrivacyMap} is parameterized as follows:
\begin{itemize}
    \item \texttt{MI} is \rustdoc{metrics/type}{L2Distance<RBig>}.
    \item \texttt{MO} is \rustdoc{measures/struct}{PrivacyCurveDP}.
\end{itemize}

\subsubsection*{Caller-verified}
None.

\subsection*{Pseudocode}
\lstinputlisting[language=Python,firstline=2,escapechar=|]{./pseudocode/NoisePrivacyMap_for_ZExpFamily2_PrivacyCurveDP.py}

\subsection*{Postcondition}
\begin{theorem}
Given a distribution \texttt{self}, the implementation returns an error if
\texttt{self} is invalid. Otherwise it returns a privacy map for the mechanism
\texttt{function(x) = x + Z}, where the coordinates of $Z$ are independent
discrete Gaussian samples represented by \texttt{self}, under
\texttt{PrivacyCurveDP}.
\end{theorem}

\begin{proof}
Line \ref{line:neg-scale} rejects a negative scale, which does not represent a
valid discrete Gaussian distribution. Line \ref{line:neg-sens} rejects a
negative distance bound.

When \texttt{d\_in} is zero, adjacent inputs are identical and line
\ref{line:zero-sens} returns the identity privacy curve. When the sensitivity
is positive and the scale is zero, the mechanism has no finite privacy
guarantee; line \ref{line:zero-scale} returns the non-private curve represented
by $\rho = \infty$.

It remains to consider positive rational sensitivity $s$ and scale $\sigma$.
By Theorem 14 of \cite{CKS20}, the discrete Gaussian mechanism is
$\rho$-zCDP for
\[
    \rho = \frac{1}{2}\left(\frac{s}{\sigma}\right)^2.
\]
Line \ref{line:rho} performs the arithmetic exactly over rationals and then
casts toward positive infinity, so the returned floating-point value
$\widehat{\rho}$ satisfies $\widehat{\rho} \geq \rho$.

Line \ref{line:curve} constructs a \texttt{PrivacyCurve} from this zCDP
guarantee. Its privacy-profile and tradeoff queries are conservative
conversions of the same zCDP guarantee. Therefore it is a valid
\texttt{PrivacyCurveDP} distance for the mechanism. Increasing the distance to
any \texttt{d\_out} accepted by the measure ordering preserves the guarantee,
which proves the privacy-map postcondition.
\end{proof}

\bibliographystyle{alpha}
\bibliography{mod}

\end{document}
